% Part: quantified-modal-logic
% Chapter: semantics
% Section: fixed-domain-models

\documentclass[../../../include/open-logic-section]{subfiles}

\begin{document}

\olfileid{qml}{sem}{fdm}

\olsection{Fixed-Domain Models}

\begin{explain}
Fixed domain semantics is the simplest way of combining the semantics
of modal logic with the semantics of first-order logic. We have a
single, fixed domain of objects (as in first-order logic) and a set of
worlds (as in modal logic). The only novelty is in the semantics of
\emph{!!{predicate}s} and \emph{!!{function}s}, which becomes
world-relative. 

In first-order logic the interpretation of a monadic !!{predicate}
meaning `is red' would be a set of objects (namely, the red ones). In
modal logic, however, we're considering multiple worlds, which may
differ as to which objects are red in them. Hence we need to associate
such a !!{predicate} to a single set of objects, but to a set of
objects \emph{for each world}. To do so, our interpretations of
!!{predicate}s are \emph{functions from worlds to} sets of (or
relations between) objects.

Similarly, a !!{function} meaning `the birthplace of' should be
associated with a one function \emph{for each world}, since which
place is the birthplace of a given individual may differ from world to
world. To do so, our interpretations of !!{function}s are
\emph{functions from worlds to} to functions from objects to objects.
\end{explain}

\begin{defn}[Fixed-domain models]
    A \emph{fixed-domain model} $\mModel M$ for the language $\Lang{L}$ 
    of quantified modal logic consists of the following elements:
    \begin{enumerate}
    \item \emph{Worlds:} a non-empty set~$\mWorlds{M}$
    \item \emph{Accessibility relation:} a binary relation~$\mAccess{M}$ on~$\mWorlds{M}$
    \item \emph{Domain:} a non-empty set, $\mDomain{M}$
\item \emph{Interpretation of !!{constant}s:} for each !!{constant}~$c$ of
  $\Lang{L}$, !!a{element} $\mAssign{c}{M} \in \mDomain{M}$
\item \emph{Interpretation of !!{predicate}s:} for each $n$-place
  !!{predicate}~$R$ of $\Lang{L}$ (other than $\eq$), a function 
  from worlds to $n$-place
  relations on the domain: $\mAssign{R}{M} \colon W \to \mDomain{M}^n$
\item \emph{Interpretation of !!{function}s:} for each $n$-place
  !!{function}~$f$ of $\Lang{L}$, an $n$-place function from worlds 
  to $n$-place functions on the domain: $\mAssign{f}{M} \colon W \to 
  (\mDomain{M}^n \to \mDomain{M})$

\end{enumerate}

We write $\mAssign{R}{M}[w]$ for the value of $\mAssign{R}{M}$ at $w$ (a $n$-place relation) and $\mAssign{f}{M}[w]$ for the value of the 
$\mAssign{f}{M}[w]$ at $w$ (a $n$-place function). 
\end{defn}

\begin{rem}[Model types]
  As in modal logic, a model $\mModel{M}$ is called a $\Ax{D}$-model, 
  $\Ax{T}$-model, etc., if its accessibility relation is serial,
  reflexive, etc., as in \olref{tab:five-types}. 
\end{rem}

\begin{table}[t]
    \begin{tabular}{| c | l | l |}
      \hline
      Model type & Property of $R$ & Characteristic schema \\
      \hline \hline
      \Ax{K} & no constraint & \\
      \hline
      \Ax{D} & \emph{serial}: $\forall u \exists v Ruv$ 
      & $\Box p \lif \Diamond p$ \\
      \hline
      \Ax{T} & \emph{reflexive}: $\forall w Rww$
      & $\Box p \lif p$ \\
      \hline
      \Ax{B} & \emph{symmetric}: $\forall u\forall v(Ruv \lif Rvu)$
      & $p \lif \Box\Diamond p$ \\
      \hline
      \Ax{4} & \emph{transitive}:
      $\forall u \forall v \forall w ((Ruv \land Rvw) \lif Ruw)$
      & $\Box p \lif \Box \Box p$ \\
      \hline
      \Ax{5} & \emph{euclidean}:
      $\forall w \forall u \forall v ((Rwu \land Rwv) \lif Ruv)$ 
      & $\Diamond p \lif \Box\Diamond p$ \\
      \hline
    \end{tabular}
    \caption{Five model types.}
    \ollabel{tab:five-types}
  \end{table} 

\end{document}